\begin{enumerate}

\item Momentos respecto a $A$: $R_B L = qL\cdot L/2 + P a$.

\item Equilibrio vertical: $R_A + R_B = qL + P$.

\item En el tramo $x > a$ el cortante es $C_y = R_A - qx - P$ y se anula en $x = 2{,}5$ m.

\item El flector máximo es el flector en ese punto.

\end{enumerate}

\noindent\textbf{Operaciones y resultados intermedios.}

\begin{enumerate}

\item Reacción en B, de los momentos respecto a A.

\[\fitmath{y_{1}=\frac{12\cdot 6\cdot 3+18\cdot 2}{6}\approx 42\,\mathrm{kN}}\]

\item Reacción en A, del equilibrio vertical.

\[\fitmath{y_{2}=12\cdot 6+18-\frac{12\cdot 6\cdot 3+18\cdot 2}{6}\approx 48\,\mathrm{kN}}\]

\item Abscisa donde se anula el cortante (x mayor que a).

\[\fitmath{y_{3}=\frac{48-18}{12}\approx 2.5\,\mathrm{m}}\]

\item Flector máximo.

\[\fitmath{y_{4}=48\cdot 2.5-\frac{12\cdot {2.5}^{2}}{2}-18\cdot \left(2.5-2\right)\approx 73.5\,\mathrm{kN m}}\]

\end{enumerate}
